\subsection{Class MultiRegisterAST}
\label{sec:MultiRegisterAST}

\declaredin{MultiRegister.h}

A \code{MultiRegisterAST} object represents an ordered collection of registers
as a single operand. As a \code{MultiRegisterAST} is an \code{Expression}, it
may contain the physical registers' contents if they are known.


\begin{apient}
typedef boost::shared_ptr<MultiRegisterAST> Ptr;
\end{apient}
\apidesc{
A type definition for a reference-counted pointer to a \code{MultiRegisterAST}.
}

\begin{apient}
MultiRegisterAST(MachRegister r, uint32_t num_elements = 1);
\end{apient}
\apidesc{
Construct a collection of \code{num\_elements} consecutive registers beginning
at \code{r}. The registers after the first are formed by incrementing the
register number of \code{r}, so a collection built this way is always
consecutive. The \code{MachRegister} datatype is Dyninst's register
representation and should not be constructed manually.
}

\begin{apient}
MultiRegisterAST(std::vector<RegisterAST::Ptr> _in);
\end{apient}
\apidesc{
Construct a collection from the registers in \code{\_in}, keeping their order.
\code{\_in} must not be empty, as its first element supplies the base register.
The registers given need not be consecutive.
}

\begin{apient}
bool isUsed(Expression::Ptr findMe) const;
\end{apient}
\apidesc{
\code{isUsed} is meant to return \code{true} if \code{findMe} is a
\code{MultiRegisterAST} holding the same registers, in the same order, as this
one. It compares the constituent \code{RegisterAST::Ptr} values rather than the
registers they refer to, so it holds only when the two objects share their
element objects. A \code{MultiRegisterAST} built
separately for the same registers returns \code{false}, which is why
\code{Operand::isRead} and \code{Operand::isWritten} do not answer \code{true}
for a candidate the caller constructed. Every other \code{Expression} in this
interface compares by value.
}

\begin{apient}
std::string format(formatStyle how = defaultStyle) const;
\end{apient}
\apidesc{
Returns the constituent registers, each formatted in turn, separated by commas
and enclosed in brackets, as in \code{[RAX,RBX]}.
}

\begin{apient}
bool operator<(const MultiRegisterAST& rhs) const;
\end{apient}
\apidesc{
Orders two \code{MultiRegisterAST} objects so that they may be placed into sets
or other sorted containers. Like \code{isUsed} it compares the constituent
\code{RegisterAST::Ptr} values rather than the registers, so the order is by
address and does not follow the register numbers.
}

\begin{apient}
RegisterAST::Ptr getBaseRegAST() const;
\end{apient}
\apidesc{
Returns the first register of the collection. For a collection constructed
from a starting register and a count, this is the register it started from.
}

\begin{apient}
uint32_t length() const;
\end{apient}
\apidesc{
Returns the number of registers in the collection.
}

\begin{apient}
const std::vector<RegisterAST::Ptr>& getRegs() const;
\end{apient}
\apidesc{
Returns the registers themselves, in the order the collection holds them.
This is the only way to reach them individually.
}

\begin{apient}
bool areConsecutive() const;
\end{apient}
\apidesc{
Returns \code{true} if the registers are known to be consecutive, which tells
a caller which of the two constructors produced this object. A collection
built from a starting register and a count is consecutive by construction;
one built from a vector of registers is not, since the registers in it are
taken as given.
}

